\documentclass[12pt]{article}
\usepackage[utf8]{inputenc}
\usepackage{amsmath}
\usepackage{geometry}
\usepackage{fancyhdr}

\geometry{letterpaper, margin=1in}
\pagestyle{empty} % No headers or footers for the first page

\begin{document}

\begin{center}
\vspace*{-0.5in} % Adjust this to move the title block up or down
{\Large \textbf{EMCH 290}}\\
{\large Mechanical Engineering Thermodynamics}\\
{\large Extra Exercise\# 1}\\
\vspace{0.2in}
{ Points Possible: 100 [2\% of FINAL Class Grade]}
\end{center}

\noindent Name: \underline{\hspace{4in}} \\
\vspace{0.3in}
\noindent Date: \underline{\hspace{2in}} 

% Include Question 1 on the first page
\vspace{-0.3in}
\subsection*{Question 1 [25 points]}
Show ALL steps in the conversion process. You are allowed to use internet conversion tables, given that sources are referenced correctly.
\begin{enumerate}
    \item Convert 3000 kJ to cal. 
    \item Convert 300 m/s to mph.
\end{enumerate}

\newpage % Starts a new page

\pagestyle{fancy} % Restores the headers and footers starting from the second page
\lhead{EMCH 290 Extra Exercise \#1}
\rhead{Page \thepage}

% Other questions appear on subsequent pages
\subsection*{Question 2 [25 points]}
Show ALL steps in the conversion process. 
\begin{enumerate}
    \item Convert 500 ft-lbf/min to W.
    \item Convert 20000 ft/s² to m/s².
\end{enumerate}
\newpage

\subsection*{Question 3 [25 points]}
On Mars, where the acceleration due to gravity is \( g = 3.72076 \, \text{m/s}^2 \), a satellite has a weight of 2000 N on Earth. Calculate the weight of the satellite on Mars.
\newpage

\subsection*{Question 4 [25 points]}
In a tank, three layers of different fluids are stacked on top of each other. The bottom layer is water (density \( = 1000 \, \text{kg/m}^3 \)), the middle layer is oil (density \( = 890 \, \text{kg/m}^3 \)), and the top layer is gasoline (density \( = 750 \, \text{kg/m}^3 \)). Each layer is 2 meters high. Calculate the pressure at the bottom of the tank, considering \( g = 9.81 \, \text{m/s}^2 \). Draw a Diagram of the tank with the densities listed appropriately.

\end{document}
